\section{Formal Properties and Adversarial Analysis}
\label{sec:formal}

The mechanism walk of \S\ref{sec:design} described what the loop does;
this section states what the loop \emph{enforces} and prices what an
adaptive adversary can still do to it. The two are paired deliberately:
each invariant names the evaluation that stresses it, and each attack
channel closes against one of the invariants.

\subsection{The Serving-Time Consistency Invariant}
\label{sec:invariant}

The fiction graph (\S\ref{sec:design}) turns G2 from an aspiration into a
data structure. We state the consistency condition it enforces, together
with its induction and its measured stress test.\begin{property}[Serving-time fiction consistency]
\label{prop:consistency}
Let $\mathcal{G}$ be a session's fiction graph and $F(\mathcal{G})$ the
set of host facts derivable from it, two facts conflicting when they assign
conflicting normalized values to one predicate. Write $\sigma_t$ for the
served response at turn $t$ and $E(\sigma_t)$ for the host facts the
extraction layer emits from that response, $E$ being the same
soundness-approximating extractor the consistency check runs (so the
invariant inherits the extractor's failure modes; see the extractor
condition below). The loop maintains the invariant
\begin{equation*}
E(\sigma_t) \not\pitchfork F(\mathcal{G}_{t-1}) \ \wedge \ F(\mathcal{G}_t) = \mathrm{cl}\big(F(\mathcal{G}_{t-1}) \cup E(\sigma_t)\big)
\end{equation*}
for every served turn, where $a \not\pitchfork b$ denotes that no fact
in $a$ conflicts with any fact in $b$ under the graph's predicate-level
normalization. The invariant is conditional on the extractor ($E$ must
surface every fact a conflict-sensitive reader would surface), so the
property is a guarantee about the checked relation, and its soundness in
the world is bounded empirically by the re-query benchmark of
\S\ref{sec:rq2}, not derived. Verification failure breaks the left
conjunct \emph{before} serving, so the violated draft is never
observable to the attacker; the failure path serves deterministic vetted
content pinned to the seed identity and never commits the rejected
draft's facts to the graph (\textsc{FailSafeFallback} in
Algorithm~\ref{alg:loop}). The argument is a short induction: the empty
graph is trivially conflict-free; if the invariant holds at $t-1$, any
draft whose emitted facts conflict with the closure is rejected before
serving, and the closure of a conflict-free set remains conflict-free;
the failure path emits only vetted facts pre-registered against the
seed graph. Because $\mathrm{cl}$ is the transitive closure of the
graph's predicate conflicts, the invariant extends across the session:
an attacker who re-queries any fact under any phrasing receives a value
consistent with every prior turn \emph{of that session} -- and strictly
of that session (\S\ref{sec:threatmodel}). The C1 benchmark
(\S\ref{sec:rq2}) tests exactly this: zero contradictions over 45
adversarial re-queries, upper confidence bound $0.079$.
\end{property}

\subsection{The Verified-Serve Algorithm}
\label{sec:algorithm}

\begin{algorithm}[t]
\caption{Verified-serve loop (per attacker turn $x$)}
\label{alg:loop}
\begin{algorithmic}[1]
\Require session state $\sigma$, fiction graph $\mathcal{G}_\sigma$, attacker turn $x$, policy $\pi$
\State $c \gets \textsc{Classify}(x, \sigma)$ \Comment{behavioral class (Chrollo)}
\State $\pi \gets \textsc{Plan}(c, \textsc{Retrieve}(x, \sigma))$ \Comment{grounded plan (Don)}
\State $B \gets \textsc{FactBlock}(\mathcal{G}_\sigma)$ \Comment{graph serialization}
\For{$\mathit{try} \in \{1, 2\}$}
  \State $d \gets \textsc{Draft}(x, \pi, B, \textsc{Steer}(\pi))$ \Comment{provider LLM}
  \State $\mathcal{G}' \gets \textsc{Observe}(\mathcal{G}_\sigma, d)$ \Comment{absorb emitted host facts}
  \If{$\textsc{DangerScan}(d) \ \vee$}
     \State \hspace{\algorithmicindent} $\textsc{IdentityPin}(d, \pi) = \bot \ \vee$
     \State \hspace{\algorithmicindent} $\textsc{Consistent}(\mathcal{G}', d) = \bot \ \vee$
     \State \hspace{\algorithmicindent} $\textsc{TerminalSanity}(d) = \bot$ \Comment{§\ref{sec:verify}}
     \State \hspace{\algorithmicindent} \textbf{continue} \Comment{retry with violation fed back}
  \EndIf
  \If{$\textsc{Require}(d, \pi)$} \Comment{high-stakes content class}
     \State \textbf{return} \textsc{HITLGate}$(d, \pi)$ \Comment{G3: human decides}
  \EndIf
  \State \textbf{return} $d$ \Comment{verified: serve}
\EndFor
\State \textbf{return} \textsc{FailSafeFallback}$(\mathcal{G}_\sigma)$
  \Comment{G4: pinned-identity vetted transcript, graph unchanged}
\end{algorithmic}
\end{algorithm}

Algorithm~\ref{alg:loop} gathers the serving path into one place: the
four decision-procedure checks of \S\ref{sec:verify} form the reject
arm, the \textsc{Require} branch carries G3, and the fallback arm
carries G4. Every branch is deterministic given the draft; the only
nondeterminism in the loop is the provider model itself, which is
exactly the component the checks exist to bound.

\subsection{Attacking the Loop}
\label{sec:attack}
An adaptive adversary (\S\ref{sec:threatmodel}) who has read this
section has two footholds, and the design prices both. The first is the
\emph{observation channel}: a served response's facts enter the graph,
so an attacker who induces the honeypot to commit a lie has poisoned one
session's world. The design contains this rather than prevents it.
Commitment is one-way on purpose: the graph absorbs what was served,
never what was merely drafted, so a rejected draft leaves no trace, and
the session lifetime bounds the poisoned world to the attacker who
created it. What the channel cannot do is leak true state, because the
graph contains no true state to leak; every triple originates in the
deception plan, the seed identity, or a prior verified serve.

The second foothold is the \emph{steering channel}. Predicted next
techniques and matched bait reach the generator as context, and
attacker-controlled input influences what is retrieved and predicted;
an attacker who knows this could shape their own session to steer bait
of their choosing. Three facts bound the exposure. Bait is drawn from a
closed canary-registered catalog, so the worst case is an attacker
obtaining a credential that was designed to be obtained and is designed
to be traced. Retrieval runs over a curated corpus, not open text, so
corpus poisoning requires access the threat model denies
(\S\ref{sec:threatmodel}). And the HITL gate sits between every
high-stakes deployment and the wire, so steering the machinery still
requires persuading a human. The residual risk is stated plainly: a
sufficiently informed attacker can make the honeypot commit deeper to
one fiction branch than the operator intended, at no safety cost and at
the cost only of that session's plan flexibility.
